\section{Sample Topic}
\label{sec:sample}
\sources{Lecture ``Sample'' (\nolinkurl{01-sample-lecture.pdf}); textbook Ch.~1.}

Organize by topic, not by lecture date; merge group slides into the topic they
cover and check them against the textbook.

\subsection{Boxes}

\begin{definition}[Term]
A definition, with a textbook citation in a footnote.\HP{\S1.1}{2}
\end{definition}

\begin{theorem}[Result]
A formula: $T=I\times\CPI\times C$.\HP[48]{\S1.9}{46}
\end{theorem}

\begin{eg}[Worked example]
Re-derive every number with a script before writing it.
\end{eg}

\begin{remark}
\slidefix{} Marks a corrected error in the slides.
\end{remark}

\beyond{} Marks material added from the textbook, not in the slides.

\subsection{Textbook Exercises}

\begin{tbexercise}{1.1}{61}{\S1.1, \hppages{}{2}{5}}
Restate the exercise here. The box cites the exercise page and the pages with the
theory needed; an appendix exercise takes the letter as the optional argument,
e.g.\ \texttt{\textbackslash begin\{tbexercise\}[B]\{B.1\}\{60\}\{...\}}.
\end{tbexercise}
